\documentclass[10pt,a4paper]{article}
\usepackage[T1]{fontenc}\usepackage{lmodern,amsmath,amssymb,amsthm,microtype}
\usepackage[margin=25mm]{geometry}
\usepackage[colorlinks=true,linkcolor=blue,urlcolor=blue]{hyperref}
\setlength{\emergencystretch}{2em}
\newtheorem{theorem}{Theorem}[section]\newtheorem{proposition}[theorem]{Proposition}
\newtheorem{corollary}[theorem]{Corollary}
\newcommand{\Z}{\mathbb Z}\newcommand{\Q}{\mathbb Q}
\title{Reference-point elimination and label-dependent progressions}
\author{}\date{14 November 2026}
\begin{document}\maketitle
\begin{abstract}
Let $P$ be a non-torsion point on an elliptic curve. For infinitely many
integer labels $n$, allow any infinite set of reference indices $h$,
with the set allowed to depend on $n$. The tuples
$([h]P,[hn]P,[(h-1)n]P)$ are Zariski dense in the third power of the
curve. Thus even label-dependent congruence restrictions leave the
basic obstruction to algebraically recovering a reference point from
the two other points. A divisor intersection argument gives an
explicit bound for the number of permitted reference indices after a
fixed proper algebraic restriction, outside a finite set of labels.
The result does not preclude a complete family with finitely many
reference choices for each label and gives no absolute lower bound
for Diophantine definitions.
\end{abstract}

\section{A quantitative statement for divisors}
Work over an algebraically closed field of characteristic zero.
Let $E$ be an elliptic curve with origin $O$, and let $P\in E$ be of
infinite order. Put $L=\mathcal O_E(O)$, a line bundle of degree one,
and write
\[
 L_d=L^{\otimes d}\boxtimes L^{\otimes d}\boxtimes L^{\otimes d}
 \quad\hbox{on }E^3.
\]
Let $D$ be the zero divisor of a nonzero section of $L_d$, with $d\ge1$.
For $n\in\Z$, put
\[
 \Gamma_n=\{(Q,[n]Q):Q\in E\}\subset E^2,
 \qquad F_n=E^2\times\{[n]P\}\subset E^3.
\]

\begin{proposition}\label{prop:bound}
There are only finitely many $n$ with $F_n\subset D$.
If $F_n\not\subset D$ and $n^2+1>2d$, then
\[
 \#\{h\in\Z:([h]P,[hn]P,[n]P)\in D\}
                       \le d(1+n^2).
\]
\end{proposition}
\begin{proof}
A fiber $E^2\times\{T\}$ contained in $D$ is an irreducible component
of that divisor, since both have dimension two. There are finitely
many such components, and $[n]P$ are distinct, proving the first claim.

For a fiber not contained in $D$, restriction gives an effective
Cartier divisor $D_n$ on $E^2$ with line bundle
$L^{\otimes d}\boxtimes L^{\otimes d}$. Use the ample line bundle
$A=L\boxtimes L$ on $E^2$. The two fiber classes have self-intersection
zero and intersect each other in one, so
\[
                          A\cdot D_n=2d.
\]
On $\Gamma_n$, the projections to $E$ pull $L$ back to line bundles
of degrees $1$ and $n^2$. Here $\deg([n]^*L)=n^2$ in characteristic
zero, with degree zero when $n=0$ and the pullback is trivial.
Consequently
\[
                        \deg(A|_{\Gamma_n})=1+n^2.
\]
If $\Gamma_n$ were a component of $D_n$, positivity of intersection
with $A$ on every nonzero effective curve would give $2d\ge1+n^2$.
The stated inequality excludes this.

The restricted section on $\Gamma_n\simeq E$ is therefore nonzero.
Its zero divisor has degree $d(1+n^2)$, so it has at most this many
geometric points. Distinct integers $h$ give distinct points $[h]P$,
which proves the bound, without any condition on the sign or
coprimality of $h,n$.
\end{proof}

\section{Infinite choices for each label}
\begin{theorem}\label{thm:dense}
Let $I\subset\Z$ be infinite. For each $n\in I$, let
$H_n\subset\Z$ be infinite, with no uniformity hypothesis on the sets
$H_n$. Then
\[
 \{([h]P,[hn]P,[(h-1)n]P):n\in I,\ h\in H_n\}
\]
is Zariski dense in $E^3$.
\end{theorem}
\begin{proof}
The map
\[
 (Q_0,Q_1,Q_2)\longmapsto(Q_0,Q_1,Q_1-Q_2)
\]
is an algebraic automorphism of $E^3$. It sends the displayed tuples
to $([h]P,[hn]P,[n]P)$. It suffices to prove density of these images.

Suppose a proper closed subset contains them. It is contained in the
zero divisor of a nonzero section of some $L_d$. One way to see this
is to embed each $E$ by $L^{\otimes3}$ and use the Segre embedding of
the product. A homogeneous polynomial vanishing on the proper closed
subset but not on $E^3$ restricts to a section of $L_{3k}$ for some
positive $k$.

By Proposition~\ref{prop:bound}, outside finitely many $n$ the fiber
is not contained in this divisor, and outside another finite set
$n^2+1>2d$. For any remaining $n\in I$, that proposition allows only
finitely many $h$, contrary to the hypothesis on $H_n$.
\end{proof}

\begin{corollary}\label{cor:crt}
For each $n$ in an infinite set, choose any integers $h_n,m_n$ with
$m_n>0$. The tuples in Theorem~\ref{thm:dense} remain dense when $h$
is restricted to $h_n+m_n\Z$, or to any infinite subset of that
progression. Finite-place congruences and additional real restrictions
which still leave infinitely many $h$ for each label do not change
this conclusion.
\end{corollary}
\begin{proof}
These are permitted choices of the infinite sets $H_n$.
\end{proof}

\section{The precise implication for a variable reduction}
A rational formula recovering the first point from the last two on
a fixed nonempty open subset would put the tuples in the graph of a
rational map $E^2\dashrightarrow E$, whose closure has dimension at
most two. A polynomial recovery rule with only finitely many choices
for the first point, on such an open subset, likewise lies in a
proper algebraic subset. Neither can hold on any of the dense families
of Theorem~\ref{thm:dense}. The preimage of that open subset is itself
dense and cannot be avoided by the displayed tuples.

More generally, impose any fixed proper algebraic restriction on
$(H,U,A)\in E^3$. After containing it in a divisor as above, all but
finitely many labels $A=[n]P$ permit at most $d(1+n^2)$ of the reference
indices $h$ used in the existing family $H=[h]P$, $U=[hn]P$.
This is stronger than saying that one particular formula failed: a
whole infinite congruence progression cannot survive such a restriction
for arbitrarily large labels.

The conclusion is nevertheless limited. A new completeness proof may
select a finite, possibly nonuniform set of reference indices for each
input. The theorem says nothing against that possibility. Nor does it
rule out replacing the source equations by a fundamentally different
construction. It does not prove that twenty-five or thirteen total
variables is minimal, and it is not a justification for stopping the
search for smaller definitions.

\section{Source relation and checking scope}
The result is an algebraic theorem and uses no prime-pattern,
contact-height, pole-realization, or $abc$ hypothesis. For its project
application, the assertion that the fixed point $P$ is non-torsion is
one of the explicitly retained elliptic inputs. The theorem refers to
the existing family of witnesses, not to every rational zero of the
project's polynomials.

The exact checker accompanying
\nolinkurl{joint_pole_realization_20261114_v29.tex}
checks the group-coordinate change on rational examples, including
negative and zero indices and identity points. These examples are
consistency checks only. The proof for all labels is the divisor
argument above, not a finite point search.
\end{document}
